\section{Učení bez učitele}\label{sec:ai-unsupervised}

Při učení bez učitele známe pouze vstupní objekty
\[
    \mathbf{x}_1,\ldots,\mathbf{x}_n\in\R^p,
\]
ale nemáme k~nim připojený výstupní soubor \(\mathbf{Y}\). Cílem proto není napodobit známé správné odpovědi, nýbrž objevit užitečnou strukturu dat. Může jít o~shluky podobných objektů, odlehlá pozorování, menší počet skrytých znaků nebo často se opakující vztahy.

Na rozdíl od učení s~učitelem zde neexistuje jediná předem známá správná odpověď, s~níž bychom mohli výsledek jednoduše porovnat. Dvě různá rozdělení mohou být smysluplná pro dva různé účely: zákazníky lze například seskupovat podle nákupního chování, podle věku nebo podle místa bydliště. Před interpretací výsledku proto musíme rozumět použitým znakům a~otázce, kterou datům klademe.

\begin{symboldef}
    Pro vektor \(\mathbf{u}=(u_1,\ldots,u_p)\in\R^p\) budeme symbolem
    \[
        \norm{\mathbf{u}}
        =
        \sqrt{\sum_{r=1}^{p}u_r^2}
    \]
    označovat jeho \emph{eukleidovskou normu}, tedy jeho délku. Vzdálenost vektorů \(\mathbf{u}\) a~\(\mathbf{v}\) je potom \(\norm{\mathbf{u}-\mathbf{v}}\). Dvojité svislé čáry tedy v~této kapitole vždy vyjadřují délku vektoru, nikoli absolutní hodnotu čísla.
\end{symboldef}

\subsection{Shlukování a~algoritmus \(k\)-means}

Shlukování rozděluje objekty do skupin tak, aby si objekty uvnitř stejného shluku byly podobné a~objekty z~různých shluků odlišné. Podobnost obvykle vyjadřujeme vzdáleností. Výsledek proto vždy závisí na zvoleném popisu objektů, jejich měřítku a~konkrétní vzdálenosti.

Algoritmus \(k\)-means hledá \(k\) shluků
\[
    C_1,\ldots,C_k
\]
a~jejich středy
\[
    \boldsymbol{\mu}_1,\ldots,\boldsymbol{\mu}_k.
\]
Přiřazení objektu \(\mathbf{x}_i\) budeme značit \(c(i)\in\set{1,\ldots,k}\). Shluk číslo \(j\) je tedy množina
\[
    C_j
    =
    \set{\mathbf{x}_i:c(i)=j}.
\]
Je-li \(C_j\neq\emptyset\), jeho střed definujeme po složkách jako aritmetický průměr všech vektorů ve shluku:
\[
    \boldsymbol{\mu}_j
    =
    \frac{1}{\sizeof{C_j}}
    \sum_{\mathbf{x}_i\in C_j}\mathbf{x}_i.
\]
Například první souřadnice \(\boldsymbol{\mu}_j\) je průměrem prvních souřadnic objektů z~\(C_j\). Pokud během výpočtu některý shluk zůstane prázdný, průměr není definován; jednoduchá implementace proto ponechá jeho dosavadní střed, případně lze zvolit nový střed z~dat.

Pro eukleidovskou vzdálenost probíhá opakování dvou kroků:
\begin{enumerate}
    \item každý objekt přiřadíme k~nejbližšímu středu;
    \item každý střed nahradíme aritmetickým průměrem objektů přiřazeného shluku.
\end{enumerate}

\begin{algorithm}
    \caption{\textsc{k-means}}
    \label{alg:ai-kmeans}
    \KwIn{Objekty \(\mathbf{x}_1,\ldots,\mathbf{x}_n\in\R^p\), počet shluků \(k\) a~počáteční středy \(\boldsymbol{\mu}_1,\ldots,\boldsymbol{\mu}_k\)}
    \While{\textup{některé přiřazení nebo střed se změní}}{
        \ForEach{\textup{objekt \(\mathbf{x}_i\)}}{
            \(c(i)\gets\argmin_{j\in\set{1,\ldots,k}}\norm{\mathbf{x}_i-\boldsymbol{\mu}_j}\)
        }
        \For{\(j\gets1,2,\ldots,k\)}{
            \(C_j\gets\set{\mathbf{x}_i:c(i)=j}\)\\
            \If{\(C_j\neq\emptyset\)}{
                \(\displaystyle\boldsymbol{\mu}_j\gets\frac{1}{\sizeof{C_j}}\sum_{\mathbf{x}_i\in C_j}\mathbf{x}_i\)
            }
        }
    }
    \KwOut{Shluky \(C_1,\ldots,C_k\) a~jejich středy \(\boldsymbol{\mu}_1,\ldots,\boldsymbol{\mu}_k\)}
\end{algorithm}

\begin{figure}[ht]
    \centering
    \begin{tikzpicture}[every node/.style={font=\normalsize}]
    \node[task,fill=blue!10,minimum width=34mm] (start) at (0,0)
        {počáteční středy\\\(\boldsymbol{\mu}_1,\ldots,\boldsymbol{\mu}_k\)};
    \node[task,fill=orange!15,minimum width=38mm] (assign) at (5,0)
        {přiřazení bodů\\k~nejbližšímu středu};
    \node[task,fill=green!12,minimum width=38mm] (update) at (10,0)
        {nové středy\\jako průměry};
    \draw[diredge] (start) -- (assign);
    \draw[diredge] (assign) -- (update);
    \draw[diredge,rounded corners=5mm]
        (update.south) -- ++(0,-1) -- node[below]{opakování} ++(-5,0) -- (assign.south);
\end{tikzpicture}

    \caption{Jedna iterace algoritmu \(k\)-means.}
    \label{fig:ai-kmeans}
\end{figure}

Algoritmus se snaží zmenšovat součet čtvercových vzdáleností
\[
    J
    =
    \sum_{j=1}^{k}
    \sum_{\mathbf{x}_i\in C_j}
    \norm{\mathbf{x}_i-\boldsymbol{\mu}_j}^{2}.
\]
Pro pevně zvolené shluky je nejlepším středem jejich aritmetický průměr; jde o~vícerozměrnou obdobu
\smartref[tvrzení~][]{\showrefs}{prop:ai-prumer-vlastnosti}{minimalizační vlastnosti průměru}.
Přiřazovací krok zase pro pevné středy volí pro každý objekt nejmenší možný příspěvek do \(J\). Hodnota \(J\) se tedy v~každém kroku nezvětší. Protože existuje jen konečně mnoho přiřazení \(n\) objektů do \(k\) shluků, algoritmus se po konečném počtu změn zastaví.

\warningbox{Výsledné rozdělení nemusí být nejlepší ze všech možností. Závisí na počátečních středech, předem zvoleném počtu \(k\) a~na měřítku znaků. Algoritmus je vhodný hlavně pro přibližně kulové a~podobně velké shluky.}

\begin{remark}
    Název \(k\)-means zavedl James MacQueen roku 1967. Velmi podobný iterační postup popsal Stuart Lloyd už roku 1957 pro kvantování signálu; jeho práce však vyšla tiskem až roku 1982. Proto se algoritmus někdy označuje také jako Lloydův algoritmus.
\end{remark}

\subsection{Jednoduchá implementace}

Program~\ref{lst:ai-kmeans} zachycuje přímo oba kroky algoritmu. Pro jednoduchost pracuje s~dvourozměrnými body, výpočet vzdálenosti je však možné stejným způsobem rozšířit na libovolné \(p\).

Počáteční středy zde tvoří prvních \(k\) bodů. Taková volba je snadno čitelná, ale v~praxi může vést k~nevhodnému lokálnímu minimu. Běžnou alternativou je několik náhodných spuštění nebo inicializace \(k\)-means++, která vybírá nové středy s~ohledem na jejich vzdálenost od středů již zvolených.

\lstinputlisting[
    caption={Jednoduchá implementace algoritmu \(k\)-means.},
    label={lst:ai-kmeans}
]{07-ai/code/kmeans.py}

Funkce \texttt{assign} vyhledá nejbližší střed, zatímco \texttt{update} z~přiřazených bodů vytvoří nové průměry. Iterace skončí, jakmile se středy nezmění více než o~zadanou toleranci.

Výstupem nejsou názvy tříd, ale pouze čísla shluků. Číslo samotné nemá věcný význam: při jiném spuštění si mohou dva shluky svá čísla prohodit, aniž by se změnilo jejich rozdělení. Význam skupin proto vždy určujeme až následnou interpretací jejich středů a~typických objektů.
